% Decimal commas of the source ("0,2734") are rendered as decimal points.
\begin{description}
\item[1)] The star is a W UMa type, so both primary and secondary minima were
  observed. If the depth of the minimum isn't given, there is no way of
  deciding which (primary or secondary it is).

\item[2)] From the data, the can be seen that on four nights more than one
  minimum was observed. This allows us to estimate half the period based on
  the difference between pairs of observations:

  \begin{center}
  \begin{tabular}{ccc}
    \hline
    minima & $\Delta T$ & average \\
    \hline
    2--1   & .1367 & $\Sigma/5 = .1367$ \\
    4--3   & .1372 & Period $P =$ \\
    10--9  & .1364 & obtained average \\
    11--10 & .1373 & $\times\,2 = \underline{0.2734}$ \\
    12--11 & .1359 & \\
    \hline
  \end{tabular}
  \end{center}

\item[3)] For a first approximation of the elements, we chose an initial epoch
  ($M_0$), for example the first observed minimum 2454092.4111. The
  approximate elements are then:
  \[ M_{\mathrm{calc}} = 2454092.4111 + 0.2734 \cdot E \]

\item[4)] For an exact determination of the the elements we draw an O--C
  diagram. The data for making the O--C diagram are calculates as follows:

  \begin{center}
  \begin{tabular}{cccccc}
    \hline
    No & $(M_{\mathrm{obs}})$ & $E = (M_{\mathrm{obs}} - M_0)/P$
       & $E$ (rounded & $M_{\mathrm{calc}}$
       & O--C $=$ \\
       & $2454000+$ & & to the nearest & $2454000.0+$
       & $M_{\mathrm{obs}} - M_{\mathrm{calc}}$ \\
       & & & whole or half & & \\
       & & & integer) & & \\
    \hline
    1  & 092.4111 & 0         & 0.0    & 092.4111 & 0.0000 \\
    2  & 092.5478 & 0.5       & 0.5    & 092.5478 & 0.0000 \\
    3  & 367.3284 & 1005.5497 & 1005.5 & 367.3148 & 0.0136 \\
    4  & 367.4656 & 1006.0516 & 1006.0 & 367.4515 & 0.0141 \\
    5  & 388.5175 & 1083.0519 & 1083.0 & 388.5033 & 0.0142 \\
    6  & 388.6539 & 1083.5508 & 1083.5 & 388.6400 & 0.0139 \\
    7  & 704.8561 & 2240.1061 & 2240.0 & 704.8271 & 0.0290 \\
    8  & 776.4901 & 2502.1178 & 2502.0 & 776.4579 & 0.0322 \\
    9  & 835.2734 & 2717.1262 & 2717.0 & 835.2389 & 0.0345 \\
    10 & 847.3039 & 2761.1295 & 2761.0 & 847.2685 & 0.0354 \\
    11 & 847.4412 & 2761.6316 & 2761.5 & 847.4052 & 0.0360 \\
    12 & 847.5771 & 2762.1287 & 2762.0 & 847.5419 & 0.0352 \\
    \hline
  \end{tabular}
  \end{center}

\item[5)] The O--C diagram should be plotted on graph paper.
  % FIGURE NEEDED: O-C diagram, (O-C) in days vs epoch E, rising roughly
  % linearly from 0 at E = 0 to ~0.0355 d at E = 2762 (to be drawn from the
  % table above; 2011 DA solutions, p. 82, scan).

\item[6)] From the diagram, it can be seen that over 2762 epochs the value of
  (O--C) increased by 0.0355 days, therefore over one cycle (one epoch) the
  period was too short by:
  \[ 0.0355/2762 = 0.0000129 \]
  and the estimated period $P$ should be increased by this value. Taking into
  account that the graph can be read to 3 significant figures, the correction
  should also be to 3 significant figures.
  \[ P' = 0.2734 + 0.0000129 = 0.2734129 \]
  thus the corrected light elements of V1107 Cas are:
  \[ M_{\mathrm{calc}} = 2454092.4111 + 0.2734129 \cdot E \]

\item[7)] The correction of the period could also be determined directly from
  the table, using the mean value of O--C for the initial epoch ($E = 0$; 0.5)
  and the final epochs ($E = 2761$; 2761.5; 2762).

\item[8)] To calculate the ephemeris, we need to know the JD of a given day.
  Since dates and JD numbers are given in the observations, we can calculate
  that, for example on 1 January 2008 at 12:00 the JD was 2454467.0, while on
  1 January 2011 at 12:00 it was 2455563.0. Since 1 September 2011 is the
  243-rd day of the year, thus at 12:00 on that day JD 2455806.0 begins.

\item[9)] On 1 September 2011 at 18:00 (JD 2455806.25) there have been
  1713.8389 days or 6268.3176 epochs since the inital epoch $M_0$
  (2454092.4111). Thus the nearest (secondary) minimum which agrees with
  these elements will be:
  \[ \mathrm{JD}\ 2454092.4111 + 0.2734129 \cdot 6268.5 = 2455806.2999
     = \underline{19^{\mathrm{h}}\,12^{\mathrm{m}}\ \mathrm{UT}} \]
  The next (primary) minimum will be:
  \[ \mathrm{JD}\ 2454092.4111 + 0.2734129 \cdot 6269 = 2455806.4366
     = \underline{22^{\mathrm{h}}\,29^{\mathrm{m}}\ \mathrm{UT}} \]
  The next (secondary) minimum will be on the same night, at:
  \[ \mathrm{JD}\ 2455806.4111 + 0.2734129 \cdot 6269.5 = 2457520.5733 \]
  % sic — so printed in the source; the first term should read
  % 2454092.4111, giving JD 2455806.5733.
  that is on 2 September at
  $\underline{1^{\mathrm{h}}\,46^{\mathrm{m}}\ \mathrm{UT}}$.

  Earlier and later minima can be determined by adding or subtracting half a
  period, or $3^{\mathrm{h}}17^{\mathrm{m}}$.
\end{description}
